% ============================================================================
\section*{Resumen ejecutivo}
\addcontentsline{toc}{section}{Resumen ejecutivo}
\markboth{Resumen ejecutivo}{}
% ============================================================================

Cada sesión, el sistema hace cinco cosas en un orden fijo. Primero \emph{mide} lo que descuentan
las opciones: reconstruye, a partir de cotizaciones de compra y venta, la distribución implícita
$\Q$ del rendimiento a un plazo constante, con bandas que reflejan los spreads. Después
\emph{describe el cambio} respecto de la sesión anterior y separa la parte que ya explicaban el
precio, la volatilidad y el mercado. Luego \emph{pronostica} la distribución física $\P$ de los
rendimientos posteriores al instante de decisión y comprueba si la información de opciones la
mejora. Con ello \emph{decide}, comparando mantener la posición actual con modificarla bajo
incertidumbre, costos y riesgo conjunto de la cartera. Por último \emph{registra} todo lo necesario
para reproducir la corrida.

\begin{ideaclave}[La decisión arquitectónica central]
Mantener separadas tres capas: la \textbf{medición $\Q$}, la \textbf{predicción $\P$} y la
\textbf{acción sobre la posición}. $\Q$ entra a los modelos como variable, nunca como probabilidad
física por decreto; $\P$ se estima con rendimientos realizados y se calibra fuera de muestra; la
acción pasa siempre por costos, riesgo y una regla explícita de abstención.
\end{ideaclave}

\paragraph{Hipótesis que se quiere contrastar} Las \emph{innovaciones} en la forma de las
distribuciones implícitas, una vez descontada la información ya observable en precio, volatilidad y
mercado, pueden mejorar los pronósticos de riesgo o de rendimiento. Es una hipótesis que hay que
contrastar, no una ventaja demostrada: el documento describe cómo medirla y qué resultados la
refutarían. La primera pregunta que conviene responder es si la innovación de las colas, medida con
incertidumbre y relativa al mercado, mejora la gestión de una posición existente.

\paragraph{Qué contiene y qué no} Una arquitectura, un conjunto de métodos y un protocolo de
evidencia. No contiene resultados predictivos ni modelos entrenados. Las gráficas se construyen con un
mundo sintético de parámetros conocidos (apéndice~\ref{ap:reproducibilidad}); su función es mostrar
qué calcula cada etapa y qué puede salir mal, no anticipar el desempeño.

\paragraph{Estado al 28 de septiembre de 2026} Desde que se escribió esta propuesta se construyó la captura
diaria con datos reales de Alpaca, y quedó instalada y verificada en un repositorio privado. La
\secref{sec:estado} resume qué está hecho, qué cambió respecto de la propuesta y qué falta. Todavía no hay
resultados sobre la capacidad predictiva.

\paragraph{Cómo leer este documento} Cada sección corresponde a una etapa del proceso y empieza con
una ficha de \emph{entradas, proceso, salidas y controles}. Cinco tipos de recuadro se repiten:

\begin{center}\small
\begin{tabularx}{\linewidth}{@{}l>{\RaggedRight\arraybackslash}X@{}}
\iconoidea\ \textsf{\textbf{Idea clave}} & el concepto que conviene retener de la etapa.\\
\iconoalerta\ \textsf{\textbf{Qué no concluir}} & interpretaciones tentadoras que el método no permite.\\
\iconocheck\ \textsf{\textbf{Criterio para avanzar}} & condición verificable antes de pasar a la siguiente etapa.\\
\iconopendiente\ \textsf{\textbf{Pendiente de definir}} & decisiones del usuario que parametrizan el sistema.\\
\icononota\ \textsf{\textbf{Nota técnica}} & detalle de implementación o advertencia matemática.\\
\end{tabularx}
\end{center}

La \figref{fig:mapa} resume el recorrido completo. Los colores se mantienen en todos los diagramas:
azul para la medición $\Q$, naranja para la predicción $\P$ y verde para la acción.

\begin{figure}[H]
\centering
\resizebox{\linewidth}{!}{% Mapa del proceso: medición Q, predicción P y acción, con capas transversales.
\begin{tikzpicture}[
  n/.style={caja=#1,text width=2.28cm,minimum height=11.5mm},
  etq/.style={cajallena=#1,text width=1.62cm,minimum height=11.5mm,font=\sffamily\tiny\bfseries},
  node distance=4.2mm and 4.2mm]
% --- Medición Q -------------------------------------------------------------
\node[etq=qazul] (lq) {MEDICIÓN\\ Q};
\node[n=qazul,right=of lq] (q1) {Snapshots bid/ask\\ sincronizados};
\node[n=qazul,right=of q1] (q2) {Calidad, forward\\ y desamericanización};
\node[n=qazul,right=of q2] (q3) {Superficie sin\\ arbitraje (SSVI)};
\node[n=qazul,right=of q3] (q4) {Densidad, cuantiles\\ y bandas};
\node[n=qazul,right=of q4] (q5) {Transporte, factores\\ e innovación};
\foreach \a/\b in {q1/q2,q2/q3,q3/q4,q4/q5} \draw[flecha] (\a) -- (\b);
% --- Predicción P -----------------------------------------------------------
\node[etq=qnaranja,below=10mm of lq] (lp) {PREDICCIÓN\\ P};
\node[n=qnaranja,right=of lp] (p1) {Información base:\\ precio, vol., mercado};
\node[n=qnaranja,right=of p1] (p2) {Modelos P por\\ objetivo y horizonte};
\node[n=qnaranja,right=of p2] (p3) {Combinación y\\ calibración temporal};
\node[n=qnaranja,right=of p3] (p4) {Pronósticos P\\ calibrados};
\node[n=qnaranja,right=of p4] (p5) {Regímenes y\\ confianza medible};
\foreach \a/\b in {p1/p2,p2/p3,p3/p4,p4/p5} \draw[flecha] (\a) -- (\b);
% --- Acción -----------------------------------------------------------------
\node[etq=qaqua,below=10mm of lp] (la) {ACCIÓN};
\node[n=qaqua,right=of la] (a1) {Escenarios\\ conjuntos P};
\node[n=qaqua,right=of a1] (a2) {Costo, riesgo y\\ rendimiento (CVaR)};
\node[n=qaqua,right=of a2] (a3) {Zona de no\\ operación};
\node[n=qaqua,right=of a3] (a4) {Recomendación\\ o abstención};
\node[n=qaqua,right=of a4] (a5) {Paper, registro y\\ reconciliación};
\foreach \a/\b in {a1/a2,a2/a3,a3/a4,a4/a5} \draw[flecha] (\a) -- (\b);
% --- Conexiones entre capas -------------------------------------------------
\draw[flecha] (q5.south) -- ++(0,-0.32) -| (p2.north);
\draw[flecha] (p4.south) -- ++(0,-0.32) -| (a1.north);
\draw[flecha] (p5.south) -- ++(0,-0.6) -| (a4.north);
% --- Capas transversales ----------------------------------------------------
\node[caja=tenue,below=7mm of a1.south west,anchor=north west,minimum width=14.85cm,text width=14.5cm,
  minimum height=7mm] (reg) {\textbf{Registro reproducible}: hashes de entrada, versión de modelo, parámetros,
  hora de corte y motivo de abstención en cada corrida};
\node[caja=tenue,below=2mm of reg,minimum width=14.85cm,text width=14.5cm,minimum height=7mm] (conf)
  {\textbf{Confianza}: calidad de datos $\cdot$ sensibilidad del ajuste $\cdot$ incertidumbre predictiva
  $\cdot$ estabilidad de la decisión};
\end{tikzpicture}
}
\caption{Mapa del proceso diario. Las flechas entre capas son deliberadamente estrechas: de $\Q$ a
$\P$ solo pasan variables candidatas, y de $\P$ a la acción solo pasan distribuciones físicas
calibradas. El registro y la evaluación de confianza atraviesan todas las etapas.}
\label{fig:mapa}
\end{figure}
